\documentclass{article}

\usepackage[dblblindworkshop]{neurips_2026}
\workshoptitle{Rethinking AI for and from the Global South (GlobalSouthAI)}

\usepackage[utf8]{inputenc}
\usepackage[T1]{fontenc}
\usepackage{hyperref}
\usepackage{url}
\usepackage{amsmath}
\usepackage{microtype}

% Tighten the title block for a one-page abstract; rules and fonts unchanged.
\makeatletter
\renewcommand{\@toptitlebar}{\hrule height 4\p@ \vskip 0.14in \vskip -\parskip}
\renewcommand{\@bottomtitlebar}{\vskip 0.15in \vskip -\parskip \hrule height 1\p@ \vskip 0.05in}
\makeatother

\title{Which Component Carries the Effect?\\
Ablating an ITS-Inspired Generation Framework}

\author{Anonymous Author(s)}

\begin{document}
\maketitle
\vspace{-2.7em}

\noindent
Design Science Research: we design and evaluate a reusable framework for generating
curriculum-constrained learning material at class and learner level. Inspired by
Intelligent Tutoring System (ITS) component models, it separates a \textbf{curriculum
model}, a \textbf{lesson model}, an \textbf{expert model} and a \textbf{learner model},
and combines them through a generation-and-validation process. We instantiate it on
Nigerian junior-secondary (JSS1) mathematics and add one component at a time. We are aware
of no published component-additive ablation of an educational LLM generation system; it
lets us ask which component carries the effect.

\vspace{-0.78em}
\section{Framework, instantiation, evaluation}
\vspace{-0.52em}
The framework's intended applicability is generic; the evidence below is specific to the
evaluated instantiation.

Each rung adds one component: \textbf{C0} curriculum model, objectives only;
\textbf{C1} $+$ lesson model, the teacher's lesson plan and its instructional
specification; \textbf{C2} $+$ expert model, scoped retrieval over her lesson note, the
subject knowledge that realises the plan; \textbf{C3} $+$ a validation-and-repair agent,
a processing component, not a further knowledge model. All rungs share one
generator, artefact spec and output contract, so a rung's effect is the component's, not a
change of system. A second part adds the learner model:
\textbf{P1} personalises frozen class-level material (\textbf{P0}) from a pupil's marked
script --- score, solution steps, errors, misconceptions.

\vspace{-0.78em}
\section{Evidence from the prototype}
\vspace{-0.52em}
Nigerian JSS1 mathematics: forty pupils, one teacher. Six lessons $\times$ 4 conditions
$\times$ 3 draws $=$ 72 packs over 17 objectives, at
temperature 0 with caching disabled. Primary: the lesson-plan-based family
(C1--C3) against the objectives-only baseline.

A blinded LLM judge, run twice independently, is repeat-stable (movement 0.076;
92.4\% of 288 scores identical). The family beats C0 by $+0.338$, positive in all
six lessons (Wilcoxon $p=.031$). The component diagnostics locate it: \textbf{the lesson
model carries all of it} --- C1$-$C0 $=+0.368$ (6/6 lessons, $p=.031$), while
C2$-$C1 $=-0.042$ (2/6) and C3$-$C2 $=-0.007$ (1/6). Sharpest as a floor: 8 of 18 C0 packs
carry a judged score below 4, against \textbf{0 of 18} at each of C1, C2 and C3. Component execution is
verified without a rater: source refs on 168 C2 and 157 C3 blocks against none at C0/C1,
and a validator logging 4.22 findings per pack that closed its loop on only 6 of 18.

Four practising mathematics teachers, blinded and counterbalanced (ANOVA on position $F=0.01$, $p=.998$), returned 1{,}199 of 1{,}224 scores. On them every contrast
is \textbf{null}, including the pre-registered C1$-$C0 $=+0.044$ (95\% CI
$[-0.195,+0.283]$). That panel could not have resolved the contrast either way: rater
identity carries 25.9\% of score variance against the packet's 14.6\%,
$\mathrm{ICC}(1,1)=-0.163$, 91\% of scores are a 4 or a 5, and the minimum detectable
effect ($\approx$0.34) exceeds the largest present (0.152). We report both instruments, not the one that separates.

\vspace{0.15em}
\noindent\textbf{Contributions.} (i) A reusable ITS-inspired framework separating curriculum, lesson, expert and learner
models behind one generation-and-validation process, with the component-additive ablation
that makes each separately measurable; (ii) a prototype built from artefacts a Nigerian
school already produces, with corpus and blinded protocol; and (iii) the observation that
none of COGENT [1], FACET [2] or Lect\=uraAgents [3] reports inter-rater reliability or
validates its judge --- a property the field assumes. The learner-model rung is instantiated and under evaluation.

{\scriptsize\setlength{\parindent}{0pt}\vspace{0.1em}
\textbf{References.}
[1]~Liu et al., \emph{BEA/ACL}, 2025.\quad
[2]~Gonnermann-M\"uller et al., \emph{arXiv:2508.11401}.\quad
[3]~Sesay et al., \emph{arXiv:2606.16428}.
\par}

\end{document}
